% Standalone section -- paste into the weekly report, or \input it.
% Compressed summary of bifurcation_volume.tex for the weekly narrative: the one number and the
% one figure, without the dominance stratification, the angle comparison, or the remodeling
% discussion. Read bifurcation_volume.tex for the full working note. Numbers and figure from
% scripts/bifurcation_volume.py and scripts/plot_volume_example.py, run 2026-09-16 (800 cases).
% Second half added the same day: the same frustum volume, read over each named artery's whole
% length rather than a bifurcation-local window, from scripts/named_segment_volume.py.

\section*{Local branch volume differs between diseased and healthy arteries}

At the crux of the heart (CRUX), where the posterior descending and posterolateral arteries
branch, the main daughter's lumen volume is larger in diseased cases than in healthy ones:
median 7.78\,mm$^3$ against 6.74\,mm$^3$ (+15.5\%), Mann--Whitney $U$, $q = 0.003$ after
Benjamini--Hochberg correction across all 18 bifurcation-by-field tests run together
($n_h = 400$, $n_d = 371$). Two of the other four named bifurcations separate the same way
after correction; the full result, including why the effect reverses sign at one of them, is in
bifurcation\_volume.tex.

Volume is read off the same 3\,mm window of centerline, past the bifurcation core, that
bifurcation angle and caliber are already measured from. Between each pair of consecutive
centerline points in that window, the vessel is treated as a frustum (a cone with the tip cut
off) of the two points' radii, and the frustum volumes are summed along the window.

\Cref{fig:week3-volume-example} shows two real branches at that volume: one healthy case and one
diseased case, each the closest match in the cohort to its own group's median.

\begin{figure}[htbp]
  \centering
  \includegraphics[width=\linewidth]{content/methods/topology/figures/volume_example.png}
  \caption{CRUX main-daughter lumen radius over the same window the volume is integrated across,
    for one real healthy case (case 335, 6.75\,mm$^3$) and one real diseased case (case 54,
    7.78\,mm$^3$), each the closest match in the cohort to its group's median volume. The shaded
    area is the quantity the Mann--Whitney test above compares.}
  \label{fig:week3-volume-example}
\end{figure}

\subsection*{The effect holds over the whole vessel, not just the junction}

If disease enlarges these arteries rather than the bifurcation window happening to land on a
locally affected spot, the same separation should show up over each vessel's whole length, not
only the 3\,mm read near CRUX. Repeating the same frustum-volume sum over every point of each of
the 13 named coronary arteries (\texttt{docs\_thesis/label\_map.json}, labels 1--13), rather than
a bifurcation-local window, confirms this: 6 of the 26 name-by-field tests (Mann--Whitney $U$,
Benjamini--Hochberg FDR across all 26 together) survive correction, and every one of them is in
the distal right-coronary/crux territory, R-PDA, R-PLA, RCA and L-PLA, the same two vessels CRUX
itself joins (\Cref{tab:week3-vessel-volume}). None of the proximal left arteries (LM, LAD, LCX)
separate at all. The whole R-PDA and R-PLA are larger in diseased cases, not just the stretch
nearest the junction, so the CRUX result above is the near-junction edge of a whole-vessel effect,
not something specific to the bifurcation itself (\texttt{scripts/named\_segment\_volume.py}).

\begin{table}[h]
  \centering
  \caption{Whole-vessel volume and radius, diseased vs.\ healthy, for the 6 of 26 name
    $\times$ field tests (13 named arteries $\times$ \{volume, radius\}) surviving
    Benjamini--Hochberg FDR (\texttt{scripts/named\_segment\_volume.py}).}
  \label{tab:week3-vessel-volume}
  \begin{tabular}{llrrrrr}
    \hline
    Artery & Field & $n_h$ & $n_d$ & healthy & diseased & $q$ \\
    \hline
    R-PDA & volume & 379 & 365 & 74.9\,mm$^3$  & 92.1\,mm$^3$  & 0.003 \\
    R-PDA & radius & 379 & 365 & 0.673\,mm     & 0.697\,mm     & 0.003 \\
    R-PLA & volume & 378 & 358 & 104.6\,mm$^3$ & 127.4\,mm$^3$ & 0.003 \\
    R-PLA & radius & 378 & 358 & 0.724\,mm     & 0.750\,mm     & 0.003 \\
    RCA   & radius & 412 & 386 & 1.187\,mm     & 1.244\,mm     & 0.008 \\
    L-PLA & volume & 31  & 39  & 27.6\,mm$^3$  & 68.4\,mm$^3$  & 0.023 \\
    \hline
  \end{tabular}
\end{table}

L-PLA's separation rests on far fewer cases than the rest of this table (31 healthy, 39 diseased,
against 358+ for R-PDA/R-PLA), since L-PLA only exists in left- and co-dominant hearts; it is
reported for completeness, not weighted the same as the other five rows.
